\documentclass[12pt,a4paper]{article}

% ─── Packages ───────────────────────────────────────────────────────────────
\usepackage[utf8]{inputenc}
\usepackage[T1]{fontenc}
\usepackage[french]{babel}
\usepackage{geometry}
\usepackage{amsmath}
\usepackage{amssymb}
\usepackage{booktabs}
\usepackage{array}
\usepackage{xcolor}
\usepackage{colortbl}
\usepackage{hyperref}
\usepackage{tcolorbox}
\usepackage{enumitem}
\usepackage{titlesec}
\usepackage{fancyhdr}
\usepackage{graphicx}
\usepackage{caption}
\usepackage{subcaption}
\usepackage{float}
\usepackage{pgfplots}
\usepackage{tikz}
\usepackage{multirow}
\usepackage{tabularx}
\usepackage{lmodern}

\pgfplotsset{compat=1.18}

% ─── Page geometry ──────────────────────────────────────────────────────────
\geometry{
  top=2.5cm, bottom=2.5cm,
  left=2.5cm, right=2.5cm,
  headheight=15pt
}

% ─── Colors ─────────────────────────────────────────────────────────────────
\definecolor{primaryblue}{RGB}{30, 100, 200}
\definecolor{accentorange}{RGB}{230, 100, 20}
\definecolor{lightgray}{RGB}{245, 245, 245}
\definecolor{darkgray}{RGB}{80, 80, 80}
\definecolor{successgreen}{RGB}{40, 160, 60}
\definecolor{warningyellow}{RGB}{255, 193, 7}
\definecolor{infoblue}{RGB}{173, 216, 230}
\definecolor{rowblue}{RGB}{220, 235, 255}

% ─── Section formatting ─────────────────────────────────────────────────────
\titleformat{\section}
  {\large\bfseries\color{primaryblue}}
  {\thesection}{1em}{}[\titlerule]

\titleformat{\subsection}
  {\normalsize\bfseries\color{accentorange}}
  {\thesubsection}{1em}{}

\titleformat{\subsubsection}
  {\normalsize\bfseries\color{darkgray}}
  {\thesubsubsection}{1em}{}

% ─── tcolorbox styles ───────────────────────────────────────────────────────
\tcbuselibrary{skins, breakable}

\newtcolorbox{objectifbox}{
  colback=blue!5!white, colframe=primaryblue,
  fonttitle=\bfseries, title=Objectif final,
  rounded corners, boxrule=1pt
}

\newtcolorbox{notebox}{
  colback=yellow!10!white, colframe=warningyellow,
  fonttitle=\bfseries, title=Note sur les données,
  rounded corners, boxrule=1pt
}

\newtcolorbox{formulebox}{
  colback=gray!8!white, colframe=darkgray,
  rounded corners, boxrule=0.8pt
}

\newtcolorbox{conclusionbox}{
  colback=blue!8!white, colframe=primaryblue,
  fonttitle=\bfseries,
  rounded corners, boxrule=1pt
}

\newtcolorbox{actionbox}{
  colback=orange!10!white, colframe=accentorange,
  fonttitle=\bfseries, title=Action recommandée,
  rounded corners, boxrule=1pt
}

\newtcolorbox{risquebox}{
  colback=red!5!white, colframe=red!60!black,
  fonttitle=\bfseries, title=Gestion des Risques,
  rounded corners, boxrule=1pt
}

\newtcolorbox{termbox}[1]{
  colback=blue!10!white, colframe=primaryblue!70!black,
  fonttitle=\bfseries, title=#1,
  rounded corners, boxrule=1pt
}

% ─── Header / Footer ────────────────────────────────────────────────────────
\pagestyle{fancy}
\fancyhf{}
\fancyhead[L]{\small\textcolor{darkgray}{TimeSeries Project}}
\fancyhead[R]{\small\textcolor{darkgray}{Stratégie Gigafactories --- Marché EV Japonais}}
\fancyfoot[C]{\thepage}
\renewcommand{\headrulewidth}{0.4pt}

% ─── Hyperref setup ─────────────────────────────────────────────────────────
\hypersetup{
  colorlinks=true,
  linkcolor=primaryblue,
  urlcolor=primaryblue,
  citecolor=primaryblue,
  pdftitle={Stratégie de Construction de Gigafactories de Batteries},
  pdfauthor={JDAIDI OUSSAMA}
}

% ════════════════════════════════════════════════════════════════════════════
\begin{document}

% ─── Title page ─────────────────────────────────────────────────────────────
\begin{titlepage}
  \centering
  \vspace*{1cm}

  {\Large\bfseries Université Moulay Ismaïl \hfill École Nationale Supérieure d'Arts et Métiers}

  \vspace{1.5cm}
  \rule{\textwidth}{2pt}
  \vspace{0.4cm}

  {\Huge\bfseries Stratégie de Construction\\[0.3em]
  de Gigafactories de Batteries}

  \vspace{0.4cm}
  \rule{\textwidth}{2pt}

  \vspace{0.5cm}
  {\large Marché EV Japonais $\bullet$ Prévision 2025--2035}

  \vspace{2cm}
  \begin{minipage}{0.45\textwidth}
    \centering
    \textbf{Réalisé par :}\\[0.3em]
    {\Large JDAIDI OUSSAMA}
  \end{minipage}
  \hfill
  \begin{minipage}{0.45\textwidth}
    \centering
    \textbf{Encadré par :}\\[0.3em]
    {\Large TAWFIK MASROUR}
  \end{minipage}

  \vspace{2.5cm}
  \begin{tabular}{ll}
    \textbf{Méthodes :}  & Séries Temporelles \& ML/DL \\
    \textbf{Modèles :}   & SARIMA, ARIMAX, Ridge, LSTM \\
    \textbf{Comparés :}  & 10 modèles \\[0.6em]
    \textbf{Sources :}   & JAMA $\bullet$ US EPA Fuel Economy \\
    \textbf{Période :}   & 2010 -- 2035 \\
    \textbf{Date :}      & Mai 2025 \\
  \end{tabular}
\end{titlepage}

% ─── Table of contents ──────────────────────────────────────────────────────
\tableofcontents
\newpage

% ════════════════════════════════════════════════════════════════════════════
\section{Idée \& Objectif du Projet}

\subsection{Contexte}

La transition énergétique mondiale vers les véhicules électriques crée une demande croissante et
structurelle en batteries lithium-ion. Le Japon, berceau des constructeurs automobiles mondiaux
(Toyota, Honda, Nissan), représente un marché stratégique : ses OEMs produisent plus de
\textbf{18 millions de véhicules par an} à l'échelle mondiale, soit un ratio overseas/domestique de
$4:1$. Une décision de construction de gigafactory au Japon n'alimente donc pas seulement le
marché local, mais les chaînes d'assemblage mondiales.

\begin{objectifbox}
Développer une stratégie quantifiée pour la construction de gigafactories de batteries :
\textbf{combien d'usines construire}, quand les lancer (délai de construction $\approx 3$ ans),
et dans quel scénario de demande (base / optimiste / pessimiste), en se basant sur des
prédictions de la demande batterie jusqu'en 2035.
\end{objectifbox}

\subsection{Formule centrale}

La logique du projet repose sur une équation simple mais puissante :

\begin{formulebox}
\begin{equation}
  \frac{\text{Demande GWh}(t)}{40\ \text{GWh/usine}} = N_{\text{usines}}(t)
\end{equation}
\end{formulebox}

\subsection{Problématiques traitées}

\begin{itemize}
  \item \textbf{Qui sont les clients du marché japonais ?} Le marché est domestique à 93--95\,\%
        (marques japonaises). Mais la vraie question est la demande mondiale des OEMs japonais.

  \item \textbf{Y a-t-il un risque de saturation ?} Les ventes BEV ont chuté de $-32{,}5\,\%$
        en 2024 --- signal de saturation domestique. Le marché interne seul ne justifie pas
        plusieurs usines.

  \item \textbf{Quelle est la position du Japon dans le marché global ?} Production overseas
        $4\times$ la production domestique. Les usines doivent alimenter les lignes mondiales.

  \item \textbf{Quels facteurs externes influencent la demande ?} 7 facteurs : démographie,
        politique EV, BYD, taux de change, réseau de recharge, solid-state, risque matières
        premières.
\end{itemize}

\begin{notebox}
\textbf{Limitation importante identifiée :} \texttt{vehicles.csv} est une base de données
EPA américaine utilisée comme proxy pour les capacités de batteries des modèles japonais.
Les données de ventes quotidiennes sont des interpolations de données mensuelles JAMA.
Ces limitations sont documentées et prises en compte dans l'interprétation des résultats.
\end{notebox}

% ════════════════════════════════════════════════════════════════════════════
\section{Pipeline Complet du Projet}

\begin{table}[H]
\centering
\caption{Pipeline complet du projet}
\begin{tabular}{clll}
\toprule
\textbf{\#} & \textbf{Étape} & \textbf{Statut} & \textbf{Outputs} \\
\midrule
01 & Collecte Raw Data       & \textcolor{successgreen}{\checkmark} Fait
   & \texttt{vehicles.csv}, \texttt{japan\_car\_sales}, \texttt{raw\_material\_prices} \\
02 & Nettoyage Data          & \textcolor{successgreen}{\checkmark} Fait
   & \texttt{dataCleaning.py} $\to$ fichiers clean \\
03 & Calcul Battery Demand   & \textcolor{successgreen}{\checkmark} Fait
   & \texttt{battery\_demand.csv} \\
04 & Facteurs Externes       & \textcolor{successgreen}{\checkmark} Fait
   & \texttt{battery\_demand\_enriched.csv} \\
05 & Analyse Marché Japonais & \textcolor{successgreen}{\checkmark} Fait
   & EDA $\cdot$ saisonnalité $\cdot$ saturation $\cdot$ exports \\
06 & Stationnarité \& SARIMA & \textcolor{successgreen}{\checkmark} Fait
   & forecast 2025--2030 \\
07 & ARIMAX Enrichi          & \textcolor{successgreen}{\checkmark} Fait
   & impact facteurs externes \\
08 & ML/DL XGBoost·Prophet   & \textcolor{successgreen}{\checkmark} Fait
   & \texttt{TimeSeries\_Forecasting\_Analysis.ipynb} \\
09 & Scénarios \& Décision   & \textcolor{successgreen}{\checkmark} Fait
   & Base $\cdot$ Optimiste $\cdot$ Pessimiste $\to$ $N_{\text{usines}}$/année \\
10 & Rapport Final           & \textcolor{successgreen}{\checkmark} Fait
   & Stratégie $\cdot$ Timing $\cdot$ Budget risque \\
\bottomrule
\end{tabular}
\end{table}

% ════════════════════════════════════════════════════════════════════════════
\section{Phase 1 --- Sources \& Nettoyage des Données}

\subsection{Sources de données}

\begin{table}[H]
\centering
\caption{Sources de données utilisées}
\begin{tabular}{lllllp{5cm}}
\toprule
\textbf{Fichier} & \textbf{Source} & \textbf{Lignes} & \textbf{Col} & \textbf{Période} & \textbf{Rôle} \\
\midrule
\texttt{vehicles.csv}
  & US EPA FE & 49\,846 & 84 & 2010--2024
  & Capacité batterie par modèle EV \\
\texttt{japan\_car\_sales\_2010\_2024.csv}
  & JAMA & 5\,073 & 5 & 2010--2024
  & Ventes BEV / PHEV / HEV / Total \\
\texttt{raw\_material\_prices\_2008\_2024.csv}
  & Marché mondiaux & 816 & 6 & 2008--2024
  & Japon Prix Li $\cdot$ Co $\cdot$ Ni $\cdot$ Graphite \\
\bottomrule
\end{tabular}
\end{table}

\subsection{Nettoyage des données}

\subsubsection{\texttt{car\_models\_clean.csv}}
\begin{itemize}
  \item Filtrage types EV (BEV/PHEV/HEV) depuis \texttt{vehicles.csv}
  \item Calcul $\texttt{battery\_capacity\_kwh} = \texttt{charge240} \times 7{,}2\ \text{kW}$
  \item Extraction \texttt{motor\_power\_kw} depuis la colonne texte \texttt{evMotor}
  \item Simplification \texttt{vehicle\_class} (SUV / Pickup / Sedan / \ldots)
  \item Suppression lignes sans \texttt{battery\_capacity} ni \texttt{electric\_range}
\end{itemize}

\subsubsection{\texttt{japan\_car\_sales\_clean.csv}}
\begin{itemize}
  \item Parse de la colonne \texttt{Date} $\to$ \texttt{datetime}
  \item Tri chronologique + ajout colonnes \texttt{year} / \texttt{month}
  \item Interpolation linéaire des valeurs manquantes ; Clip(\texttt{lower=0})
  \item Cast en entier pour BEV / PHEV / HEV / Total\_LDV
\end{itemize}

\subsubsection{\texttt{raw\_material\_prices\_clean.csv}}
\begin{itemize}
  \item Standardisation des noms matériaux (Title Case)
  \item Interpolation par groupe matériau (linéaire)
  \item Remplissage résiduel par médiane ; Suppression prix $\leq 0$ ; Arrondi à 2 décimales
\end{itemize}

% ════════════════════════════════════════════════════════════════════════════
\section{Calcul de la Demande Batterie}

La demande batterie est calculée en deux temps :

\begin{formulebox}
\begin{align}
  \bar{C}(t,\,\text{type}) &= \overline{\texttt{battery\_capacity\_kwh}}
    \text{ groupé par [year, type]} \tag{1} \\[4pt]
  E_{\text{BEV}}(t)  &= V_{\text{BEV}}(t)  \times \bar{C}_{\text{BEV}}(t)  \tag{2} \\
  E_{\text{PHEV}}(t) &= V_{\text{PHEV}}(t) \times \bar{C}_{\text{PHEV}}(t) \tag{3} \\
  E_{\text{HEV}}(t)  &= V_{\text{HEV}}(t)  \times \bar{C}_{\text{HEV}}(t)  \tag{4} \\[4pt]
  E_{\text{total}}(t) &= \bigl[E_{\text{BEV}}(t)
    + E_{\text{PHEV}}(t) + E_{\text{HEV}}(t)\bigr] \div 10^{6} \tag{5}
\end{align}
\end{formulebox}

% ════════════════════════════════════════════════════════════════════════════
\section{Facteurs Externes}

\begin{table}[H]
\centering
\caption{Facteurs externes et leurs formules d'application}
\begin{tabular}{lp{5.5cm}ccc}
\toprule
\textbf{Variable} & \textbf{Formule} & \textbf{BEV} & \textbf{PHEV} & \textbf{HEV} \\
\midrule
\texttt{demo\_index}
  & $(1-0{,}006)^{\max(0,\,t-2012)}$
  & $\times 1{,}0$ & $\times 1{,}0$ & $\times 1{,}0$ \\
\texttt{policy\_score}
  & $\sigma\!\bigl(0{,}6\,(t-2020)\bigr)$
  & $\times 1{,}0$ & $\times 0{,}6$ & $\times 0{,}2$ \\
\texttt{byd\_pressure}
  & $0{,}25\,\sigma\!\bigl(1{,}5\,(t-2022)\bigr)$
  & $\times 1{,}0$ & $\times 0{,}3$ & --- \\
\texttt{charger\_density}
  & $C_0\,e^{0{,}166(t-2010)}/150\,000$
  & $\times 1{,}0$ & $\times 0{,}5$ & --- \\
\texttt{solid\_state\_proba}
  & $\sigma\!\bigl(1{,}0\,(t-2027)\bigr)$ [scénario haussier]
  & $+40\%$ & $+20\%$ & --- \\
\texttt{material\_risk}
  & $0{,}6\,\sigma_{\text{Li}} + 0{,}4\,\sigma_{\text{Co}}$ (3 ans)
  & --- & --- & --- \\
\texttt{jpy\_usd}
  & Données historiques 2010--2024
  & coût ¥ & coût ¥ & coût ¥ \\
\bottomrule
\end{tabular}
\end{table}

\subsection{Formule de demande ajustée}

\begin{formulebox}
\begin{align}
D_{\text{BEV\_adj}}(t)  &= E_{\text{BEV}}(t) \times \text{demo}
  \times (1 + 1{,}0\,\text{policy}) \times (1 - 1{,}0\,\text{byd})
  \times (1 + 1{,}0\,\text{charger}) \div 10^6 \tag{6} \\[4pt]
D_{\text{PHEV\_adj}}(t) &= E_{\text{PHEV}}(t) \times \text{demo}
  \times (1 + 0{,}6\,\text{policy}) \times (1 - 0{,}3\,\text{byd})
  \times (1 + 0{,}5\,\text{charger}) \div 10^6 \tag{7} \\[4pt]
D_{\text{HEV\_adj}}(t)  &= E_{\text{HEV}}(t)  \times \text{demo}
  \times (1 + 0{,}2\,\text{policy}) \div 10^6 \tag{8} \\[4pt]
D_{\text{total\_adj}}(t) &= D_{\text{BEV\_adj}} + D_{\text{PHEV\_adj}} + D_{\text{HEV\_adj}} \tag{9} \\[4pt]
D_{\text{BEV\_ss}}(t)   &= D_{\text{BEV\_adj}}(t)
  \times \bigl(1 + 0{,}40 \times \text{solid\_state\_proba}(t)\bigr) \tag{10}
\end{align}
\end{formulebox}

% ════════════════════════════════════════════════════════════════════════════
\section{Résultats Clés de la Phase 1}

\subsection{Marché japonais --- Chiffres clés 2024}

\begin{table}[H]
\centering
\begin{tabular}{cccc}
\toprule
\textbf{1,89 M} & \textbf{154k} & \textbf{$-32{,}5\,\%$} & \textbf{14\,\%} \\
Véhicules LDV & Ventes BEV & Variation BEV 2024 & Saturation marché BEV \\
\midrule
\textbf{31,3 GWh} & \textbf{48,2 GWh} & \textbf{1,20} & \textbf{7,26} \\
Prédiction 2025 & Prédiction 2026 & Usines 2026 & Usines 2030 \\
\bottomrule
\end{tabular}
\end{table}

\subsection{Enseignements sur le marché japonais}

\begin{itemize}
  \item \textbf{Saturation domestique réelle :} Le marché BEV domestique n'est saturé qu'à
        14\,\% du potentiel ($K = 1{,}11$\,M unités/an) mais la chute de $-32{,}5\,\%$ en 2024
        signale un retournement de court terme.

  \item \textbf{La vraie variable à prédire :} La demande domestique (19,4 GWh/2024)
        n'est que 20\,\% de la demande mondiale des OEMs japonais. Il faut multiplier par le
        ratio overseas/domestique $\approx 4$.
\end{itemize}

% ════════════════════════════════════════════════════════════════════════════
\section{Analyse Statistique (EDA)}

\subsection{Tests de stationnarité}

\begin{table}[H]
\centering
\caption{Tests de stationnarité sur la série \texttt{Dadj\_GWh} journalière}
\begin{tabular}{lccc}
\toprule
\textbf{Test} & \textbf{Statistique} & \textbf{p-value} & \textbf{Conclusion} \\
\midrule
ADF (Augmented Dickey-Fuller)
  &  0,3432 & 0,9792 & \textcolor{red}{Non-Stationnaire} \\
KPSS
  &  6,3782 & 0,0100 & \textcolor{red}{Non-Stationnaire} \\
Phillips-Perron
  & $-17{,}1892$ & 0,0000 & \textcolor{successgreen}{Stationnaire} \\
\bottomrule
\end{tabular}
\end{table}

\begin{formulebox}
ADF et KPSS concordent : la série est \textbf{non-stationnaire}. Phillips-Perron donne un
résultat contradictoire dû à la rupture structurelle de 2022.
\textbf{Décision :} log-transformation + différenciation $d = 1$.
\end{formulebox}

\textbf{Observations sur la série (Figure 1 : demande journalière ajustée BEV 2010--2024) :}
\begin{itemize}
  \item Croissance stable (2010--2021)
  \item Explosion exponentielle (2022--2024), clairement visible sur le graphique
\end{itemize}

\textbf{Facteurs structurels identifiés :}
\begin{itemize}
  \item \textbf{BYD change la donne :} BYD entre en 2022 avec des prix 20--30\,\% inférieurs
        et des specs supérieures (autonomie 300\,km vs 180\,km pour le Nissan Sakura).
  \item \textbf{Solid-state = rupture :} Toyota approuvé 2024, déploiement 2027.
        $+40\,\%$ de capacité kWh par BEV si déployé.
  \item \textbf{Yen faible = coûts élevés :} $79 \to 151$\,¥/\$ entre 2012 et 2024.
\end{itemize}

\subsection{ACF \& PACF --- Analyse visuelle}

L'ACF sur la série brute décroît très lentement ($\approx 0{,}8$--$0{,}9$ à lag\,=\,40),
signature d'une série avec tendance. La PACF montre un pic dominant au lag\,=\,1 et lag\,=\,7.

% ════════════════════════════════════════════════════════════════════════════
\section{Méthodes Statistiques Classiques}

\subsection{SARIMA --- Modèle baseline}

\begin{formulebox}
\centering
\begin{tabular}{cccc}
$(0,1,0)$ & $(1,1,1)[12]$ & $-61{,}68$ & $59{,}3\,\%$ \\
Ordre AR/I/MA & Ordre saisonnier (P,D,Q) & AIC & MAPE test \\
\end{tabular}
\end{formulebox}

\subsection{ARIMAX --- Modèle enrichi avec facteurs externes}

L'ARIMAX reprend les mêmes ordres SARIMA en ajoutant deux variables exogènes
standardisées : \texttt{policy\_score} et \texttt{charger\_lag12}.

\begin{formulebox}
\centering
\begin{tabular}{cccc}
$-726$ & $72{,}6\,\%$ & \texttt{policy+charger} & 2026 \\
AIC ARIMAX & MAPE test & Variables exogènes & Seuil 1 usine (SARIMA médian) \\
\end{tabular}
\end{formulebox}

\subsection{Comparaison SARIMA vs ARIMAX}

\begin{table}[H]
\centering
\caption{Comparaison SARIMA vs ARIMAX}
\begin{tabular}{lccc}
\toprule
\textbf{Critère} & \textbf{SARIMA(0,1,0)(1,1,1)[12]}
  & \textbf{ARIMAX + policy + charger} & \textbf{Gagnant} \\
\midrule
AIC (train)        & $-61{,}68$ & $-72{,}60$ & \textcolor{successgreen}{ARIMAX} \\
MAE test (GWh)     & 0,721  & 0,901  & \textcolor{primaryblue}{SARIMA} \\
RMSE test (GWh)    & 0,820  & 1,040  & \textcolor{primaryblue}{SARIMA} \\
MAPE test (\%)     & 59,3\% & 72,6\% & \textcolor{primaryblue}{SARIMA} \\
Forecast 5 ans     & Stable & Explose & \textcolor{primaryblue}{SARIMA} \\
Interprétabilité   & Limitée & Riche  & \textcolor{successgreen}{ARIMAX} \\
Usage recommandé   & Forecast 2025--2030 & Analyse des facteurs & --- \\
\bottomrule
\end{tabular}
\end{table}

% ════════════════════════════════════════════════════════════════════════════
\section{Machine Learning}

\subsection{Approche générale ML}

\begin{table}[H]
\centering
\caption{Features utilisées dans les modèles ML}
\begin{tabular}{lll}
\toprule
\textbf{Feature} & \textbf{Optimisé par Optuna ?} & \textbf{Description} \\
\midrule
\texttt{lag\_1} \ldots\ \texttt{lag\_k}
  & Oui ($k \in [1, 30]$) & Valeurs passées de \texttt{Dadj\_GWh} \\
\texttt{rolling\_mean\_w}
  & Oui ($w \in [2, 30]$) & Moyenne des $k$ dernières valeurs \\
\texttt{policy\_score}  & Non (Phase 1) & Score politique EV (sigmoid) \\
\texttt{charger\_density} & Non (Phase 1) & Densité bornes de recharge \\
\texttt{byd\_pressure}  & Non (Phase 1) & Part de marché BYD \\
\bottomrule
\end{tabular}
\end{table}

\subsection{Décomposition Tendance-Résidus}

\begin{formulebox}
\begin{enumerate}
  \item \textbf{Transformation logarithmique :}
        $\log_y = \log(\texttt{Dadj\_GWh} + \varepsilon)$
  \item \textbf{Tendance linéaire} sur $\log_y$ (365 derniers jours) :
        $\text{tendance}(t) = \text{pente} \times t + \text{ordonnée}$
  \item \textbf{Résidus stationnaires :}
        $\text{résidus}(t) = \log_y(t) - \text{tendance}(t)$
  \item \textbf{Entraînement sur les résidus} uniquement (bornés, stationnaires)
  \item \textbf{Reconstruction :}
        $\text{GWh} = \exp(\text{tendance\_future} + \text{résidus\_prédits})$
\end{enumerate}
\textbf{Pente de tendance :} $0{,}000259$ log-unités/jour ($\approx 9{,}9\,\%$/an)
\end{formulebox}

\subsection{Résultats des modèles ML}

\subsubsection{Ridge Regression --- Meilleur modèle}

\begin{conclusionbox}
Ridge est le \textbf{modèle recommandé} pour la prévision en production.
$R^2 = 0{,}8289$~; MAPE $= 18{,}55\,\%$.
\end{conclusionbox}

% ════════════════════════════════════════════════════════════════════════════
\section{Comparaison Complète des Modèles}

\begin{table}[H]
\centering
\caption{Classement complet des 10 modèles sur le test set 2022--2024}
\begin{tabular}{clccccc}
\toprule
\textbf{Rang} & \textbf{Modèle} & \textbf{Type}
  & $\boldsymbol{R^2}$ & \textbf{RMSE} & \textbf{MAE} & \textbf{MAPE} \\
\midrule
\rowcolor{rowblue}
1  & Ridge          & ML & 0,8289  & 0,009317 & 0,006621 & \textbf{18,55\,\%} \\
2  & SVR            & ML & 0,7197  & 0,011925 & 0,008453 & 23,33\,\% \\
3  & Reg. Linéaire  & ML & 0,6085  & 0,014091 & 0,009960 & 26,96\,\% \\
\midrule
4  & SimpleRNN      & DL & $-0{,}4451$ & 0,027075 & 0,021278 & 90,31\,\% \\
5  & Random Forest  & ML & $-1{,}4032$ & 0,034914 & 0,027964 & 64,09\,\% \\
6  & XGBoost        & ML & $-1{,}5429$ & 0,035914 & 0,029026 & 67,60\,\% \\
7  & LSTM           & DL & $-1{,}9692$ & 0,038808 & 0,031988 & 77,24\,\% \\
8  & GRU            & DL & $-2{,}3555$ & 0,041256 & 0,034794 & 88,24\,\% \\
9  & ANN            & DL & $-2{,}5486$ & 0,042426 & 0,035411 & 87,12\,\% \\
10 & CNN            & DL & $-2{,}6627$ & 0,043103 & 0,036482 & 92,72\,\% \\
\bottomrule
\end{tabular}
\end{table}

\begin{conclusionbox}
\textbf{Conclusion clé :} Les modèles ML (Ridge, SVR, Régression Linéaire) surpassent
significativement les modèles DL. La raison principale : la période d'entraînement
(2010--2021) présente une demande très faible et stable, tandis que la période de test
(2022--2024) montre une croissance explosive --- une \textbf{rupture structurelle} que
les modèles DL ne peuvent pas généraliser.
\end{conclusionbox}

% ════════════════════════════════════════════════════════════════════════════
\section{Stratégie Industrielle --- Combien d'Usines ?}

\subsection{Formule}

\begin{formulebox}
\begin{equation}
  N_{\text{usines}}(\text{année}) =
    \frac{\text{Demande\_Annuelle\_GWh}(\text{année})}{40\ \text{GWh/usine/an}}
\end{equation}
\end{formulebox}

\subsection{Analyse en trois scénarios}

\begin{table}[H]
\centering
\caption{Scénarios Base / Optimiste / Pessimiste avec échéances de décision.
  $\bigstar$ = fenêtres de décision critiques.}
\small
\begin{tabular}{cccccccl}
\toprule
\textbf{Année}
  & \textbf{Base (GWh)} & \textbf{Us. Base}
  & \textbf{Opt. (GWh)} & \textbf{Us. Opt.}
  & \textbf{Pess. (GWh)} & \textbf{Us. Pess.}
  & \textbf{Échéance} \\
\midrule
2025 & 17,1 & 0,5 & 17,1 & 0,5 & 17,1 & 0,5 & 2022 \\
2026 & 18,8 & 0,5 & 19,3 & 0,5 & 18,3 & 0,5 & 2023 \\
2027 & 20,7 & 0,6 & 21,8 & 0,6 & 19,5 & 0,5 & 2024 \\
2028 & 22,7 & 0,6 & 24,6 & 0,7 & 20,9 & 0,6 & 2025 \\
\rowcolor{rowblue}
2029 & 25,0 & 0,7 & 27,8 & 0,7 & 22,3 & 0,6 & 2026\,$\bigstar$ \\
\rowcolor{rowblue}
2030 & 27,5 & 0,7 & 31,3 & 0,8 & 23,8 & 0,6 & 2027\,$\bigstar$ \\
2031 & 30,2 & 0,8 & 35,4 & 0,9 & 25,4 & 0,7 & 2028 \\
2032 & 33,3 & 0,9 & 40,0 & 1,1 & 27,2 & 0,7 & 2029 \\
2033 & 36,5 & 1,0 & 45,0 & 1,2 & 29,0 & 0,8 & 2030 \\
2034 & 40,1 & 1,1 & 50,7 & 1,3 & 30,9 & 0,8 & 2031 \\
2035 & 44,1 & 1,1 & 57,3 & 1,5 & 33,1 & 0,9 & 2032 \\
\bottomrule
\end{tabular}
\end{table}

\subsection{Calendrier de Décision}

\begin{table}[H]
\centering
\caption{Calendrier de décision avec délai de construction de 3 ans}
\begin{tabular}{ccp{5cm}l}
\toprule
\textbf{Année Décision} & \textbf{Capacité pour}
  & \textbf{Action requise} & \textbf{Urgence} \\
\midrule
2022 & 2025 & Site Usine 1 \& permis initiés
  & En cours \\
\rowcolor{rowblue}
2026 & 2029 & Expansion Usine 1 OU décision Usine 2
  & \textbf{URGENT --- Maintenant} \\
2027 & 2030 & Sécuriser chaîne appro.\ (Li, Co)
  & Haute Priorité \\
2029 & 2032 & Évaluer besoin Usine 2 (scénario optimiste)
  & À surveiller \\
2031 & 2034 & Révision finale capacité objectif 2035
  & Stratégique \\
\bottomrule
\end{tabular}
\end{table}

\subsection{Recommandations Stratégiques}

\begin{termbox}{Court terme (2025--2026)}
Sécuriser le terrain et les permis pour l'Usine 1. Bloquer les contrats
d'approvisionnement en lithium aux prix actuels ($\sim$8\,000\,\$/t, en baisse par
rapport au pic de 64\,616\,\$ en 2022). Viser 40 GWh/an de capacité opérationnelle
d'ici 2028. \textbf{La décision d'investissement pour la capacité 2029 doit être prise
MAINTENANT.}
\end{termbox}

\vspace{0.6em}

\begin{termbox}{Moyen terme (2027--2030)}
Surveiller la demande réelle vs la prévision Ridge trimestriellement. Si la demande
suit le scénario optimiste ($>30$ GWh/an en 2030), accélérer l'Usine 2. Si pessimiste,
reporter le capex et renforcer la modularité de l'Usine 1.
\end{termbox}

\vspace{0.6em}

\begin{termbox}{Long terme (2031--2035)}
Re-prévoir annuellement avec le modèle Ridge mis à jour. Le point de contrôle 2033
(36,5 GWh base) déterminera si une deuxième usine complète de 40 GWh est
économiquement justifiée.
\end{termbox}

\vspace{0.6em}

\begin{risquebox}
Concevoir l'Usine 1 comme \textbf{modulaire} (40 GWh extensible à 80 GWh) plutôt que
de s'engager sur deux sites distincts. Cela divise par deux le risque à la baisse dans
le scénario pessimiste tout en préservant l'optionalité à la hausse.
\end{risquebox}

% ════════════════════════════════════════════════════════════════════════════
\section{Conclusion}

\begin{table}[H]
\centering
\caption{Points clés du projet}
\begin{tabular}{lp{9cm}}
\toprule
\textbf{Point clé} & \textbf{Détail} \\
\midrule
Décomposition Tendance-Résidus
  & Résout le problème d'extrapolation : les modèles prédisent des résidus bornés et
    stationnaires tandis que la tendance croît analytiquement. \\[4pt]
Ridge = Meilleur Modèle
  & $R^2 = 0{,}8289$ et MAPE$= 18{,}55\,\%$, surpassant les 5 modèles DL. \\[4pt]
ML $>$ DL sur ce jeu
  & Tous les modèles ML obtiennent un $R^2 > 0$ ; tous les modèles DL obtiennent un
    $R^2 < 0$, attribuable à la rupture structurelle de la demande après 2022. \\[4pt]
Prévision Ridge : $+158\,\%$
  & La demande annuelle passe de 17,1 GWh en 2025 à 44,1 GWh en 2035. Le Japon aura
    besoin de 1,1 usine (base) à 1,5 usine (optimiste) d'ici 2035. \\[4pt]
Décision critique : 2026
  & Avec un délai de construction de 3 ans, la décision pour la capacité 2029 doit être
    prise en 2026. La fenêtre est ouverte maintenant. \\
\bottomrule
\end{tabular}
\end{table}

\begin{actionbox}
Procéder à la décision d'investissement pour l'Usine 1 en \textbf{2026} pour une capacité
opérationnelle en \textbf{2029}. Surveiller la demande trimestriellement par rapport à la
prévision Ridge. Revoir le nombre d'usines annuellement à partir de 2030.
\end{actionbox}

\vspace{2em}
\begin{center}
\textit{TimeSeries Project $\bullet$ Rapport Final $\bullet$
Demande Batteries VE Japon $\bullet$ 2025--2035}
\end{center}

\end{document}
